\clearpage\SourceLeaf{206}
\selectlanguage{english}\fancyhead[L]{\small REVIEW.}\pdfbookmark[0]{Reviews.}{anzeigen206}
\begin{center}\textbf{R E V I E W S.}\end{center}
\par\Needspace{9\baselineskip}\begin{center}
\rule{0.52\linewidth}{0.3pt}\par\medskip
Göttingische gelehrte Anzeigen.\quad 1812 February 10.\par\medskip
\rule{0.52\linewidth}{0.3pt}
\end{center}

Logarithmic and circular functions, as the simplest kinds of transcendental functions, are those with which analysts have occupied themselves most. They deserved this honour both because they constantly enter into almost all mathematical investigations, theoretical and practical, and because of the almost inexhaustible wealth of interesting truths offered by their theory.\quad Far less work has so far been done on other transcendental functions which cannot be reduced to these, but must be regarded as distinct higher kinds, connected with them only in particular cases. Yet some such functions are no less fruitful in interesting relations, and therefore no less important to anyone who esteems analysis for its own sake; their frequent occurrence in various other investigations must likewise recommend them to anyone who prefers to ask first about practical use. Professor \textsc{Gauss} has for many years been engaged in investigations of such higher transcendental functions, whose far-reaching results confirm what has been said. A part of these investigations, admittedly a very small one in proportion, which may be regarded as an introduction, as it were, to a series of papers to be supplied in the future, he submitted on 30 January under the title

\begin{center}\emph{Disquisitiones generales circa seriem infinitam}\end{center}
\[
\begin{aligned}&1+\frac{\alpha.\beta}{1.\gamma}x
+\frac{\alpha.\alpha+1.\beta.\beta+1}{1.2.\gamma.\gamma+1}xx
\\[5pt]&+\frac{\alpha.\alpha+1.\alpha+2.\beta.\beta+1.\beta+2}{1.2.3.\gamma.\gamma+1.\gamma+2}x^3+\mathit{etc.}\end{aligned}
\]
\begin{center}\emph{Pars prior,}\end{center}
\SourceLeaf{207}
as a lecture to the Royal Society of Sciences; we wish to report here the main features of its contents.

Transcendental functions always have their true source, openly or in a hidden way, in the infinite. The operations of integration, summation of infinite series, expansion of infinite products, continued fractions extending to infinity, or, in general, the approach to a limit by operations continued without end according to definite laws --- this is the proper ground on which transcendental functions are generated; or, if one prefers another image, these are the proper paths by which they are reached. Several such paths usually lead to \emph{one} goal: the circumstances and the purposes one sets oneself must determine which to choose first or in preference to the others. The series that forms the subject of the present paper has a very wide generality. According as the quantities $\alpha$, $\beta$, $\gamma$ (which, together with $x$, the author distinguishes by the names first, second, third and fourth \emph{element}, while, for brevity, he denotes the entire function represented by the series by the symbols $F(\alpha,\beta,\gamma,x)$) are determined in one way or another, it represents algebraic, logarithmic, trigonometric or higher transcendental functions; and it may be asserted that scarcely any transcendental function has so far been investigated by analysts which could not be reduced to this series. A great many truths already discovered concerning such functions as have been considered can be derived from the general nature of the function represented by our series, and for this reason alone investigations of it would deserve the attention of geometers, although this is to be regarded only as a secondary matter, and opening access to new truths as the principal aim. The present paper contains only the first half of the author's general investigations, whose excessive extent made a division necessary. Here the series itself serves as the origin of transcendental functions, which, in the further course of the work, will be derived from a source of wider application and considered from a higher point of view. The former way of generating them, by its nature, requires restriction to the cases where the series converges, thus where the fourth element $x$, positive or negative,
\SourceLeaf{208}
does not exceed the value $1$: the other way of generating them will remove this restriction. Yet the former way already leads to a multitude of remarkable truths by a more convenient and, as it were, more elementary path, and for this reason the author has begun with it.

This first half of the investigations is divided into three sections, preceded by some general remarks. To give a sample of the wide applicability of the series, $23$ different series expansions of algebraic, logarithmic and trigonometric functions are first reduced to it, as are, as an example of higher transcendental functions, the coefficients of the series arising from the expansion of $(aa+bb-2ab\cos\varphi)^n$ and proceeding according to the cosines of multiples of $\varphi$; the latter are reduced in three different ways.

The \textit{first section} deals with the relations between series of the above form in which the values of one of the first three elements differ by one unit, while those of the remaining three are the same. The author calls such series \textit{series contiguae}; in German one might perhaps call them \textit{verwandte Reihen}. Thus each series $F(\alpha,\beta,\gamma,x)$ has six contiguous series beside it, namely $F(\alpha+1,\beta,\gamma,x)$, $F(\alpha-1,\beta,\gamma,x)$, $F(\alpha,\beta+1,\gamma,x)$, $F(\alpha,\beta-1,\gamma,x)$, $F(\alpha,\beta,\gamma+1,x)$, $F(\alpha,\beta,\gamma-1,x)$, and it is shown here that there is a linear equation between the first and any two of its contiguous series. Fifteen equations arise in this way. From this follows the important theorem that, if $\alpha'-\alpha$, $\alpha''-\alpha$, $\beta'-\beta$, $\beta''-\beta$, $\gamma'-\gamma$, $\gamma''-\gamma$ are integers, there is also a linear equation between $F(\alpha,\beta,\gamma,x)$, $F(\alpha',\beta',\gamma',x)$, $F(\alpha'',\beta'',\gamma'',x)$, and thus, generally speaking, the value of the third of these functions can be derived from the values of two of them. Some of the simplest or otherwise remarkable cases have also been collected separately here by the author.

The \textit{second section} gives transformations into continued fractions, namely for the quotients
\[
\frac{F(\alpha,\beta+1,\gamma+1,x)}{F(\alpha,\beta,\gamma,x)},\quad
\frac{F(\alpha,\beta+1,\gamma,x)}{F(\alpha,\beta,\gamma,x)},\quad
\frac{F(\alpha-1,\beta+1,\gamma,x)}{F(\alpha,\beta,\gamma,x)}
\]
to which three others can be reduced by the evidently permissible interchange of the first two elements. Almost all expansions into continued fractions known so far are only particular cases of these theorems.
\SourceLeaf{209}
Especially remarkable is the case in which $\beta=0$ is set in the second expansion. From it follows a theorem which we append here because of its wide applicability. The function $F(\alpha,1,\gamma.x)$, or, what amounts to the same thing, the series
\[
1+\frac{\alpha}{\gamma}x
+\frac{\alpha\,.\,\alpha+1}{\gamma\,.\,\gamma+1}xx
+\frac{\alpha\,.\,\alpha+1\,.\,\alpha+2}{\gamma\,.\,\gamma+1\,.\,\gamma+2}x^3+\mathrm{etc.}
\]
gives the continued fraction
\[
\dfrac{1}{\displaystyle 1-\dfrac{ax}{\displaystyle 1-\dfrac{bx}{\displaystyle 1-\dfrac{cx}{\displaystyle 1-\dfrac{dx}{1-\mathrm{etc.}}}}}}
\]
where the coefficients $a,b,c,d$, etc., proceed according to the following law:
\[
\begin{array}{ll}
a=\dfrac{\alpha}{\gamma}&b=\dfrac{\gamma-\alpha}{\gamma(\gamma+1)}\\[6pt]
c=\dfrac{(\alpha+1)\gamma}{(\gamma+1)(\gamma+2)}&d=\dfrac{2(\gamma+1-\alpha)}{(\gamma+2)(\gamma+3)}\\[6pt]
e=\dfrac{(\alpha+2)(\gamma+1)}{(\gamma+3)(\gamma+4)}&f=\dfrac{3(\gamma+2-\alpha)}{(\gamma+4)(\gamma+5)}\quad\mbox{etc.}
\end{array}
\]
Accordingly, for example, the power of a binomial, the series for $\log(1+x)$, $\log\dfrac{1+x}{1-x}$, for exponential quantities, for the arc in terms of its tangent or its sine, and others can be transformed into infinite continued fractions. The transformations into such fractions given in the \textit{Theoria motus corporum coelestium} also rest on this; their proofs are supplied here by the author.

By far the greatest part of the paper is occupied by the \textit{third section}, which deals with the value of the series when the fourth element is set $=1$. After first proving with geometrical rigour that the series for $x=1$ converges to a finite sum only when $\gamma-\alpha-\beta$ is a positive quantity, the author reduces this sum, or $F(\alpha,\beta,\gamma,1)$, to the expression $\dfrac{\Pi(\gamma-1)\,.\,\Pi(\gamma-\alpha-\beta-1)}{\Pi(\gamma-\alpha-1)\,.\,\Pi(\gamma-\beta-1)}$, where the symbol $\Pi$ indicates a distinct kind of transcendental function, whose generation the author bases on an infinite product. This function, of the greatest importance throughout analysis, is at bottom nothing other than \textsc{Euler}'s inexplicable function
\[
\Pi z=1\,.\,2\,.\,3\,.\,4\,.\,.\,.\,.\,.\,z
\]
\SourceLeaf{210}
but \textit{this} way of generating it, or definition, is, in the author's judgement, entirely inadmissible, since it has a clear meaning only for positive integer values of $z$. The foundation chosen by the author is generally applicable and gives a meaning just as clear for imaginary values of $z$ as for real ones; there is no danger whatever of falling into such paradoxes and contradictions as Mr.\ \textsc{Kramp} formerly encountered with his numerical factorials. These, as can easily be shown, can be reduced to the above function, but seem less suitable for admission into analysis than it, since they depend on three quantities, while it depends on only one and must nevertheless be regarded as just as general. The author wishes to see this transcendental function $\Pi z$ granted citizenship in analysis; perhaps choosing a name of its own for it would do most to promote this. The right to do so may be reserved for whoever makes the most important discoveries in the theory of this function, which is highly worthy of the efforts of geometers. The author has already collected here a considerable number of remarkable theorems concerning it, some of which are to be regarded as new. Space does not permit us to enter into their details here: we single out only this one, that the value of the integral $\int x^{\lambda-1}(1-x^\mu)^\nu\,\mathrm{d}x$ from $x=0$ to $x=1$ can easily be reduced to the function $\Pi$, and that all the relations found by \textsc{Euler}, in part laboriously, for such integrals can be derived with the greatest ease from the general properties of those functions; conversely, $\Pi z$, whenever $z$ is a rational quantity, can be represented by some such definite integrals.

\Needspace{20\baselineskip}No less remarkable is the function arising from the differentiation of $\Pi z$, likewise transcendental, or rather
\[
\frac{\mathrm{d}\log\Pi z}{\mathrm{d}z}=\frac{\mathrm{d}\Pi z}{\Pi z\,.\,\mathrm{d}z}
\]
which the author has denoted by $\Psi z$, and which would likewise deserve a special \textit{name}. Of the numerous remarkable properties of this function established in the paper, we mention here only the one that, in general, $\Psi z-\Psi 0$, when $z$ is a rational quantity, can be reduced to logarithms and circular functions; $\Psi 0$ itself, however, is the familiar number $0{,}5772156649\,.\,.\,.\,.$, taken negatively, investigated by \textsc{Euler} and others. The author gives it here, from a calculation of his own\SourceLeaf{211}%
, to 23 decimal places, the last of which differ somewhat from \textsc{Mascheroni}'s determination. Moreover, both $\Pi z$ and $\Psi z$ are connected with several remarkable integrals for definite values of the variable quantity.

To this third section is also appended a table, calculated under the supervision of Professor \textsc{Gauss} by Mr.\ \textsc{Nicolai} with the greatest care, for $\log\Pi z$ and for $\Psi z$, in which the argument $z$ proceeds through every hundredth from 0 to 1; it is clear from the theory of these functions that all other values can easily be reduced to these values of $z$.

